\chapter{Nonlinear SSM Validation Ladder}
\label{ch:bf-nonlinear-ssm-validation}

The nonlinear SSM ladder tests whether approximate nonlinear filtering can be
used inside HMC before DSGE perturbation solvers and economic boundaries enter
the picture.  It should use small synthetic models with known parameters and
independent references, not large application models.

\section{Controlled Models}
\label{sec:bf-nonlinear-controlled-models}

The first nonlinear validation suite is deliberately small.  It is designed to
test the filtering law, structural deterministic completion, and sigma-point
moment approximation before DSGE perturbation solvers, client adapters, or HMC
geometry enter the picture.  The suite contains three first-rung models.

\subsection{Model A: Affine Gaussian Structural Oracle}
\label{subsec:bf-nonlinear-model-a}

The affine oracle is a BayesFilter-local synthetic model:
\[
  x_t=(m_t,\ell_t)^\top,\qquad
  m_t=0.35m_{t-1}-0.10\ell_{t-1}+0.25\varepsilon_t,\qquad
  \ell_t=m_{t-1},
\]
\[
  y_t=m_t+\eta_t,\qquad
  \varepsilon_t\sim N(0,1),\qquad
  \eta_t\sim N(0,0.15^2).
\]
The initial law is Gaussian and the deterministic coordinate \(\ell_t\)
stores the previous \(m\).  The oracle is the exact linear Gaussian Kalman
likelihood after converting the structural law to its full-state LGSSM.  The
test must verify that SVD cubature, the historical diagnostic SVD-UKF branch,
the promoted strict-SPD principal-square-root UKF branch, and SVD-CUT4 collapse
to the same prediction-error likelihood and that deterministic-completion
residuals are zero up to numerical tolerance.

\subsection{Model B: Nonlinear Accumulation}
\label{subsec:bf-nonlinear-model-b}

The nonlinear accumulation model is a smooth BayesFilter-local structural
fixture:
\[
  m_t=\rho m_{t-1}+\sigma\varepsilon_t,\qquad
  k_t=\alpha k_{t-1}+\beta\tanh(m_t),
\]
\[
  y_t=m_t+k_t+\eta_t,\qquad
  \varepsilon_t\sim N(0,1),\qquad
  \eta_t\sim N(0,\sigma_\eta^2).
\]
The first V1 defaults are
\[
  \rho=0.70,\qquad \sigma=0.25,\qquad \alpha=0.55,\qquad
  \beta=0.80,\qquad \sigma_\eta=0.30.
\]
This model tests nonlinear deterministic completion without a time-indexed
transition hook.  The reference oracle is a dense one-step Gaussian
moment-projection quadrature: it integrates the augmented predictive Gaussian,
computes the transformed prediction and observation moments, and then applies
the same Gaussian update used by the sigma-point filters.  It is not an exact
nonlinear likelihood.

\subsection{Model C: Nonlinear Growth Benchmark}
\label{subsec:bf-nonlinear-model-c}

The scalar nonlinear growth benchmark is the standard model associated with
Kitagawa's Monte Carlo filtering examples \citep{kitagawa1996montecarlo} and
widely reused in sequential Monte Carlo discussions.  Its time-inhomogeneous
law is
\[
  x_t=\frac{x_{t-1}}{2}
      +\frac{25x_{t-1}}{1+x_{t-1}^2}
      +8\cos(1.2(t-1))
      +\sigma_x\varepsilon_t,
\]
\[
  y_t=\frac{x_t^2}{20}+\sigma_y\eta_t,
  \qquad \varepsilon_t,\eta_t\sim N(0,1).
\]
Because the current BayesFilter structural TensorFlow transition does not pass
an explicit time index, the V1 testing fixture uses an autonomous deterministic
phase coordinate:
\[
  \tau_t=\tau_{t-1}+1,\qquad \tau_1=1,
\]
\[
  x_t=\frac{x_{t-1}}{2}
      +\frac{25x_{t-1}}{1+x_{t-1}^2}
      +8\cos(1.2\tau_{t-1})
      +\sigma_x\varepsilon_t.
\]
The observation equation remains \(y_t=x_t^2/20+\sigma_y\eta_t\).  The first
defaults are \(x_1\sim N(0,0.2)\), \(\sigma_x=1\), and \(\sigma_y=1\).
As in Model B, the first oracle is dense one-step Gaussian moment projection,
not the exact nonlinear likelihood.

\subsection{Deferred Models}
\label{subsec:bf-nonlinear-deferred-models}

\begin{description}
  \item[Bearings-only tracking.]
    \citet{gordon1993novel} and \citet{arulampalam2002tutorial} are useful
    sources, but the test should wait until angle residual and wrapping policy
    are fixed.
  \item[Radar range-bearing tracking.]
    This should follow the bearings-only residual policy and should add
    mixed-scale observation diagnostics.
  \item[Stochastic volatility.]
    This row is admitted only after it declares the one cumulative likelihood
    target being benchmarked.  For transformed actual SV that target is the exact
    transformed likelihood built from the normalizers of
    \(
      z_t=\log(y_t^2),
      \ z_t-\log(\beta^2)-h_t\sim\log(\chi_1^2)
    \).
    Same-target direct-likelihood quadrature and Zhao--Cui comparisons are
    admissible because they approximate that scalar.  Generic sigma-point rows
    that assume additive observation covariance after \(h(x_t)\), or that replace
    the observation law by Gaussian closure, do not benchmark the intended
    actual-SV target and are therefore not acceptable stand-ins for it.
\end{description}

Deferred means that the model is useful later, not that it is rejected.  The
first V1 suite stops at Models A--C so that value, score, branch, and benchmark
tests share a stable set of fixtures.

\section{Current V1 Filter Evidence}
\label{sec:bf-nonlinear-current-v1-evidence}

The first CPU nonlinear benchmark compares three current promoted-or-reference
sigma-point value filters plus the historical diagnostic comparator:
\[
  \mathcal B=\{\text{SVD cubature},\text{historical SVD-UKF},\text{principal-square-root UKF},\text{SVD-CUT4}\}.
\]
For Model A the reference is the exact Kalman prediction-error likelihood after
structural-to-LGSSM conversion.  The benchmark rows must satisfy
\[
  \left|\ell_b(y_{1:T})-\ell_{\mathrm{KF}}(y_{1:T})\right|\approx 0,
  \qquad b\in\mathcal B,
\]
up to floating-point tolerance, and the filtered means and covariances must
match the exact Kalman recursion.  This verifies that the nonlinear
sigma-point code collapses to the correct affine Gaussian law.

For Models B and C the benchmark uses only a dense one-step Gaussian projection
reference.  If \(Q_{\mathrm{dense}}\) denotes a tensor-product Gaussian rule and
\(\Pi_b\) denotes one of the sigma-point rules, the artifact reports
\[
  \left|\ell^{(1)}_{\Pi_b}-\ell^{(1)}_{Q_{\mathrm{dense}}}\right|,\qquad
  \left\|\hat x^{(1)}_{\Pi_b}-\hat x^{(1)}_{Q_{\mathrm{dense}}}\right\|_2,
  \qquad
  \left\|\hat P^{(1)}_{\Pi_b}-\hat P^{(1)}_{Q_{\mathrm{dense}}}\right\|_F.
\]
These are approximation-quality diagnostics for the first filtering step, not
exact full-likelihood errors.  The current artifact therefore records the full
Models B--C log likelihoods as filter outputs and blocks any claim that those
values have been certified against an exact nonlinear likelihood.

The same artifact records the shared branch diagnostics
\[
  r_{\mathrm{det}},\quad r_{\mathrm{supp}},\quad
  g_{\mathrm{place}},\quad g_{\mathrm{innov}},\quad
  c_{\mathrm{floor}},\quad N_{\mathrm{points}},
\]
where \(r_{\mathrm{det}}\) is the deterministic-completion residual,
\(r_{\mathrm{supp}}\) is the residual in the inactive SVD support,
\(g_{\mathrm{place}}\) and \(g_{\mathrm{innov}}\) are placement and innovation
spectral-gap diagnostics, \(c_{\mathrm{floor}}\) counts active covariance
floors, and \(N_{\mathrm{points}}\) is the rule size.  These diagnostics keep
model-law failures separate from numerical regularization.  They are also the
gate before analytic score or HMC claims.

The analytic first-order score now has a second rung beyond the affine fixture.
For Model B we use
\[
  \theta_B=(\rho,\sigma,\beta)
\]
with \(\alpha\), the initial law, the innovation variance, and the observation
variance held fixed.  The derivative provider supplies
\[
  \frac{\partial F}{\partial x_{t-1}}
  =
  \begin{bmatrix}
    \rho & 0\\
    \beta(1-\tanh^2 m_t)\rho & \alpha
  \end{bmatrix},
  \qquad
  \frac{\partial F}{\partial \varepsilon_t}
  =
  \begin{bmatrix}
    \sigma\\
    \beta(1-\tanh^2 m_t)\sigma
  \end{bmatrix},
\]
and parameter partials
\[
  F_\rho=
  \begin{bmatrix}
    m_{t-1}\\
    \beta(1-\tanh^2 m_t)m_{t-1}
  \end{bmatrix},\quad
  F_\sigma=
  \begin{bmatrix}
    \varepsilon_t\\
    \beta(1-\tanh^2 m_t)\varepsilon_t
  \end{bmatrix},\quad
  F_\beta=
  \begin{bmatrix}
    0\\
    \tanh(m_t)
  \end{bmatrix}.
\]
The observation derivative is \(H_x=[1,1]\).  Centered finite differences of
the implemented SVD cubature, the historical diagnostic SVD-UKF path, the
promoted strict-SPD principal-square-root UKF path, and the SVD-CUT4 likelihoods
match the corresponding analytic scores on the selected smooth branch.

For Model C we use
\[
  \theta_C=(\sigma_u,\sigma_y,P_{0,x}).
\]
The derivative provider supplies
\[
  F_x =
  \begin{bmatrix}
  \frac12+25\frac{1-x^2}{(1+x^2)^2} & -9.6\sin(1.2\tau)\\
  0 & 1
  \end{bmatrix},
  \qquad
  F_\varepsilon =
  \begin{bmatrix}
    \sigma_u\\0
  \end{bmatrix},
\]
\[
  F_{\sigma_u}=(\varepsilon_t,0)',\qquad
  \partial_{\sigma_y}R=2\sigma_y,\qquad
  \partial_{P_{0,x}}P_0=\operatorname{diag}(1,0),
\]
and \(H_x=[x/10,0]\).  Score parity is certified on a nondegenerate
phase-state testing variant with positive phase variance.  The default Model C
value benchmark still has zero phase variance.  Under the smooth
simple-spectrum/no-active-floor SVD score branch, that default law remains
blocked by an active placement floor.  This is an intentional diagnostic.  The
structural fixed-support branch in Chapter~\ref{ch:bf-svd-sigma-point} instead
keeps the phase direction as a declared structural null direction: the factor
column and its derivative are zero, the pre-transition sigma-point variable is
\((x_{t-1},\varepsilon_t)\), and the propagated phase coordinate is completed
pointwise by the transition map.  On that branch, default Model C is now a
score-ready testing target for SVD cubature, SVD-UKF, and SVD-CUT4, subject to
the fixed-null covariance, fixed-null derivative, active-gap, and
finite-difference diagnostics.

Consequently, GPU/XLA and HMC evidence for nonlinear Models B--C remains
deferred until the score-ready parameter boxes and branch diagnostics have
been selected and tested.

\section{Required Outputs}
\label{sec:bf-nonlinear-required-outputs}

A nonlinear SSM HMC validation run should write:
\begin{itemize}
  \item JSON diagnostics with chain counts, warmup/draw counts, step sizes,
    acceptance rate, divergence counts, R-hat, ESS, MCSE/SD, and runtime;
  \item NPZ chain artifacts for independent audit;
  \item true-parameter coverage when synthetic truth exists;
  \item filter failure counts, conditioning-guard counts, and finite-gradient
    failures.
\end{itemize}
Medium runs are sampler-usability evidence.  Strict runs require at least four
chains, finite log targets, zero or justified divergences, R-hat below 1.01,
bulk and tail ESS above the planned thresholds, and saved artifacts.

\section{Stop Rules}
\label{sec:bf-nonlinear-stop-rules}

If the LGSSM rung fails, the nonlinear rung should not start.  If the nonlinear
SSM rung fails while LGSSM passes, the failure should be attributed first to
the approximation, nonlinear target geometry, transition hook, or derivative
path.  DSGE/NK debugging should wait until this simpler explanation is ruled
out.
